\cotefile{156-9}
\scanpage{126}
\mspage{120}
% archive sheet number: 125
% Lines 1--13; continuation from the preceding scan.
m'\deleted{\unsure{d}}\unsure{avait-il|aurait-il} trompé, pour une fois ? Il faudrait que j'y sois très attentif, chaque fois que je ferai appel à des considérations de support. Je voudrais p. ex. donner aux supports \unsure{attachés} à une figure, et leurs \unsure{relations} aux supports \unsure{attachés} à une subdivision, un ``sens'' qui ne soit affecté par le fait que les deux subdivisions ne sont ``\unsure{commensurables}'' \unsure{!!|''} --- ce ``souci'' un peu vague peut-il \deleted{\unsure{se}} préciser ?

% Lines 14--21.
On pourrait bien sûr songer à travailler plutôt avec le sous-ens. $\Sigma_{\unsure{\mathrm{ens}|\mathrm{en}}}$ de $\Sigma=\Sigma_A$, \unsure{formé} des supports \unsure{en.} constructibles, avec l'opération $\unsure{\mathcal{C}|\complement}$ (partout définie) et $\operatorname{Sup}$ et $\operatorname{Inf}$ (partiellement \unsure{définies|définie} seulement). \unsure{Cet|Peut} ens. \unsure{admet-il} une construction \unsure{indépendante} de la considé%
